% STYLE-REWRITE SAMPLE. Prose rewritten to the house style profiled from
% M. Chitre's first-author papers (see chitre_first_author/style_guide_from_papers.md).
% Every equation, symbol, label and technical claim is unchanged from
% manual_gradient_ofdm_paper_concepts.tex; only the words around them differ.
\documentclass[11pt]{article}

\usepackage[margin=0.9in]{geometry}
\usepackage[T1]{fontenc}
\usepackage[utf8]{inputenc}
\usepackage{amsmath,amssymb,amsfonts,mathtools,bm}
\usepackage{booktabs,longtable,array}
\usepackage{xcolor}
\usepackage{enumitem}
\usepackage{hyperref}

\hypersetup{colorlinks=true,linkcolor=blue!50!black,urlcolor=blue!50!black}
\setlist{leftmargin=2em}
\setlength{\parskip}{0.35em}
\setlength{\parindent}{0pt}

\newcommand{\C}{\mathbb C}
\newcommand{\R}{\mathbb R}
\newcommand{\Hh}{^{\mathsf H}}
\newcommand{\T}{^{\mathsf T}}
\newcommand{\norm}[1]{\left\lVert #1 \right\rVert}
\newcommand{\abs}[1]{\left\lvert #1 \right\rvert}
\newcommand{\diag}{\operatorname{diag}}
\newcommand{\Ree}{\operatorname{Re}}
\newcommand{\argmin}{\operatorname*{arg\,min}}

\title{Which Parts of an Underwater OFDM Receiver Should We Differentiate?\\
\large Hand-derived gradients for the coded-bit latents}
\author{Project SONIQUE}
\date{\today}

\begin{document}
\maketitle

\begin{abstract}
An underwater acoustic OFDM receiver estimates a channel, equalizes it, and
decodes a forward error correction (FEC) code.  These stages are conventionally
designed and run one after another, each using the output of the one before it.
Automatic differentiation makes it possible to write the whole chain as a single
differentiable computation, and to improve every receiver variable by gradient
descent.  However, many receiver variables are chosen rather than tuned.  These
include the number of partial-FFT views, the intercarrier interference (ICI)
span, and the sparse support of a channel.  The cancellation order is also
discrete.  Since a gradient carries no information about a discrete
choice, differentiating the entire receiver spends effort where it cannot help.
We consider which parts of the receiver are worth differentiating, and which are
better left explicit.  We derive one gradient with respect to the real latent
vector that carries the coded bits.  We show that the receivers considered here
differ only in the observation model substituted into this gradient.  Although
the hand-derived gradient covers only the coded bits and a few low-dimensional
continuous refinements, it does not cover the channel, the combiner, or the
nuisance terms.  We estimate these variables by least squares, grid search, and
proximal steps.
\end{abstract}

\section{Receiver Variables}

The receivers considered here recover coded bits from observations that depend
on the channel, receiver settings, and nuisance terms.  We write the common
problem as
\begin{equation}
  \min_{\bm z,\Theta,\bm u}
  \mathcal L(\bm z,\Theta,\bm u)
  =
  \mathcal D\big(\bm y;\bm s(\bm z),\Theta,\bm u\big)
  +\mathcal R_{\rm phys}(\Theta,\bm u)
  +\lambda_p \mathcal R_{\rm code}(\bm z),
  \label{eq:master}
\end{equation}
where $\bm z\in\R^{n_b}$ is a real latent vector for the coded bits, $\Theta$
holds the receiver physics, and $\bm u$ holds nuisance terms.  The receiver
physics includes Doppler, ICI span, channel coefficients, and partial-FFT
combiners.  The nuisance terms include impulsive noise and external interference.
The first term measures how well the receiver
explains the received signal, the second keeps the physical variables plausible,
and the third rewards bit decisions that satisfy the parity checks.

The three groups of variables behave differently under a gradient.  The latent
vector $\bm z$ is real, continuous, and high dimensional, and the loss is smooth
in it, so a gradient is informative.  Much of $\Theta$ is discrete, and the rest
enters linearly, so it is either searched or solved in closed form.  Hence we
differentiate only $\bm z$ and a small number of continuous refinements, and we
keep support selection, candidate selection, and most channel and interference
updates explicit.

\section{Coded Bits as Continuous Variables}

A parity check is a statement about bits, and a bit is not a continuous
variable.  We therefore represent each coded bit $b_i$ by a real latent $z_i$
and a relaxed bipolar value
\begin{equation}
  \xi_i=\tanh z_i,\qquad \xi_i\in(-1,1),
\end{equation}
so that a large positive or negative $z_i$ corresponds to a confident bit, and
$z_i$ near zero corresponds to a bit the receiver has not yet decided.  For
data carrier $k$ with coded-bit indices $(i_I(k),i_Q(k))$, the relaxed QPSK
symbol is
\begin{equation}
  s_k(\bm z)=\frac{1}{\sqrt 2}\big(\xi_{i_I(k)}+j\xi_{i_Q(k)}\big),
  \label{eq:qpsk}
\end{equation}
which reduces to a QPSK constellation point when both latents are large.  The
derivatives are
\begin{equation}
  \frac{\partial s_k}{\partial z_i}
  =
  \begin{cases}
    (1-\xi_i^2)/\sqrt 2, & i=i_I(k),\\
    j(1-\xi_i^2)/\sqrt 2, & i=i_Q(k),\\
    0, & \text{otherwise,}
  \end{cases}
  \label{eq:symbol_derivative}
\end{equation}
where the factor $1-\xi_i^2$ vanishes as a bit becomes confident.  A confident
bit therefore stops responding to the observation, and the gradient acts on the
bits that are still in doubt.

For a low-density parity-check (LDPC) graph with check neighbourhoods
$\mathcal N_j$, we measure code consistency by the relaxed parity residual
\begin{equation}
  \mathcal R_{\rm code}(\bm z)
  =
  \sum_j
  \left(1-\prod_{i\in\mathcal N_j}\xi_i\right)^2,
  \label{eq:parity_penalty}
\end{equation}
where the product over a check equals $+1$ when the check is satisfied by hard
bits.  The residual is therefore zero for a valid codeword and grows as the
relaxed bits drift away from one.  Let $P_j=\prod_{i\in\mathcal N_j}\xi_i$ and
$P_{j\setminus i}=\prod_{r\in\mathcal N_j\setminus i}\xi_r$.  Then
\begin{equation}
  \frac{\partial \mathcal R_{\rm code}}{\partial z_i}
  =
  -2(1-\xi_i^2)
  \sum_{j:i\in\mathcal N_j}
  (1-P_j)P_{j\setminus i},
  \label{eq:parity_gradient}
\end{equation}
where the sum runs over the checks that bit $i$ takes part in.  We do not
compute $P_{j\setminus i}$ as the ratio $P_j/\xi_i$, because a single undecided
bit makes $\xi_i$ zero and the ratio undefined.  We compute the excluded
product directly, or from prefix and suffix products, so that a zero is safe.

\section{One Gradient for Every Receiver}

The receivers considered here observe the transmitted symbols through different
physics.  However, each observation is linear once the physical variables are
fixed.  We write the observation for group $a$ as
\begin{equation}
  \widehat{\bm y}_a
  =
  \sum_{k\in\mathcal D_a}\bm a_{ak}(\Theta)s_k(\bm z)
  +\bm B_a(\Theta)\bm u,
  \qquad
  \bm e_a=\widehat{\bm y}_a-\bm y_a,
  \label{eq:local_model}
\end{equation}
where $\bm a_{ak}$ collects the response from carrier $k$ into observation $a$,
$\bm B_a$ maps the nuisance terms into the same observation, and $\bm e_a$ is
the residual.  A group is one carrier for a per-carrier receiver, and one
segment or one user for the receivers that process several at a time.

For the squared observation error
\begin{equation}
  \mathcal D_{\rm obs}=\sum_a \norm{\bm e_a}_2^2,
\end{equation}
the hand-derived gradient is
\begin{equation}
  \frac{\partial \mathcal D_{\rm obs}}{\partial z_i}
  =
  2\sum_{a:\,k(i)\in\mathcal D_a}
  \Ree\left[
    \bm e_a\Hh \bm a_{a,k(i)}
    \frac{\partial s_{k(i)}}{\partial z_i}
  \right],
  \label{eq:universal_z_gradient}
\end{equation}
where $k(i)$ is the carrier that carries bit $i$.  Only the observations that
see that carrier contribute, so the cost of the gradient grows with the number
of carriers that a receiver couples together, and not with the size of the
frame.  The gradient of the full loss follows as
\begin{equation}
  \nabla_z\mathcal L
  =
  \nabla_z\mathcal D_{\rm obs}
  +\lambda_p\nabla_z\mathcal R_{\rm code}
  +\nabla_z\mathcal R_{\rm aux},
  \label{eq:full_z_gradient}
\end{equation}
where the three terms pull the latents towards the observation, towards a valid
codeword, and towards whatever else a particular receiver requires.  We derive
\eqref{eq:universal_z_gradient} once, and each receiver below reuses it by
supplying its own $\bm a_{ak}$.

\section{What Each Receiver Contributes}

Each receiver considered here supplies an observation model and a decision that
we do not differentiate.  Table~\ref{tab:concepts} lists both.

\begin{longtable}{@{}>{\raggedright\arraybackslash}p{0.24\linewidth}>{\raggedright\arraybackslash}p{0.28\linewidth}>{\raggedright\arraybackslash}p{0.38\linewidth}@{}}
\caption{What each receiver contributes to the hand-derived gradient, and what
it decides without one.\label{tab:concepts}}\\
\toprule
Receiver & What we differentiate & What stays explicit \\
\midrule
\endfirsthead
\toprule
Receiver & What we differentiate & What stays explicit \\
\midrule
\endhead
Partial-FFT demodulation
& The coded-bit latents, through the partial-FFT observation model
$\bm Y_k\approx\sum_q \bm c_{kq}s_q(\bm z)$, using
\eqref{eq:universal_z_gradient}.
& The partial-FFT windows, the number of partial-FFT views $M$, and the combiner weights.  We
choose these from candidates, or solve them by least squares. \\

Progressive ICI equalization
& The latents at the current ICI half-bandwidth $K_{\rm ICI}$.  Soft code
constraints take the place of hard tentative decisions.
& The span $K_{\rm ICI}$ itself.  Increasing it is a model-selection step, and
we decide it from the cyclic redundancy check (CRC), belief propagation (BP)
success, parity residual, and residual energy. \\

Nonuniform Doppler compensation
& The latents, once a Doppler candidate has been chosen.  A single scalar
gradient for residual Doppler can be derived later, if the interpolation is
fixed.
& Coarse resampling, timing, and the Doppler search, which we run as a grid or
from the preamble.  We do not differentiate the acquisition path. \\

Sparse delay-Doppler channel estimation
& The latents, once a sparse support or dictionary slice has been chosen.
Reliable soft symbols then act as anchors.
& Orthogonal matching pursuit (OMP) support selection, basis-pursuit active
sets, and dictionary pruning.  We run these as OMP, BP or proximal steps, and
follow them with least squares. \\

Clustered channel adaptation
& Low-dimensional cluster delay and Doppler offsets, by
$\partial L/\partial\theta_c=2\Ree[\bm e\Hh(\partial A/\partial\theta_c)\bm s]$.
& Cluster assignment and large support changes, which we treat as discrete
candidates.  Cluster amplitudes are least squares. \\

Impulsive noise mitigation
& The latents, after the current impulse estimate has been subtracted.  A soft
impulse mask has an explicit scalar gradient.
& The impulse positions, which we detect by threshold, generalized likelihood
ratio test (GLRT), or null-tone residual.  The impulse samples are least squares
or a sparse proximal update. \\

Partial-band partial-block interference
& The latents, through
$\widehat{\bm y}=A\bm s(\bm z)+B_{\rm int}\alpha$.  The code penalty keeps the
receiver from cancelling valid signal energy along with the interference.
& The interference support in time and frequency, which comes from a GLRT or a
candidate set.  The interference coefficients $\alpha$ are least squares. \\

Nonbinary LDPC coding
& The binary-image bit latents and the binary parity gradient.  Latents can
later be grouped into GF($2^m$) symbols with a nonbinary check surrogate.
& The nonbinary BP messages and the field operations.  We keep the finite-field
algebra visible rather than inside a differentiator. \\

Asynchronous multiuser OFDM
& A separate latent vector $\bm z^{(u)}$ and parity penalty for each user.  The
gradient for user $u$ uses that user's mixing matrix $A_u(\tau_u,\Theta_u)$.
& User timing offsets, overlap segmentation, and cancellation order, which are
candidates or small grid searches. \\

Physical-layer network coding
& Separate latents $(\bm z_A,\bm z_B)$, network-coded latents $\bm z_\oplus$, or
composite latents with an exclusive-or consistency penalty.
& Which decoding mode to trust.  We score candidates rather than descend a
gradient. \\

Adaptive modulation and coding, with outer erasure coding
& Nothing at packet level.  We use the receiver outputs as features: parity
residual, log-likelihood ratio margin, selected $M$, selected $K_{\rm ICI}$, CRC
status, and residual energy.
& Mode selection and erasure-code allocation, which are system-layer policies. \\
\bottomrule
\end{longtable}

\section{Losses for Four Receivers}

\subsection{Partial FFT and ICI}

Under residual Doppler, a carrier is observed together with its neighbours, and
the partial-FFT views resolve that mixture in time.  Let $\bm Y_k\in\C^M$ be the
partial-FFT vector for carrier $k$, and let $\bm c_{kq}\in\C^M$ be the response
from transmitted carrier $q$ to observed carrier $k$.  With
$\mathcal B_K(k)=\{q:\abs{q-k}\le K\}$, the residual on carrier $k$ is
\begin{equation}
  \bm r_k
  =
  \sum_{q\in\mathcal B_K(k)}\bm c_{kq}s_q(\bm z)-\bm Y_k,
\end{equation}
where the sum runs over the neighbours included in the ICI at carrier
$k$.  We minimize
\begin{equation}
  \mathcal L_{\rm pfft}
  =
  \lambda_d\sum_k\norm{\bm r_k}_2^2
  +\lambda_p\mathcal R_{\rm code}(\bm z),
\end{equation}
where $\lambda_d$ and $\lambda_p$ set how much the receiver trusts the
observation against the parity checks.  For bit $i$ on carrier $q$,
\begin{equation}
  \frac{\partial \mathcal L_{\rm pfft}}{\partial z_i}
  =
  2\lambda_d
  \sum_{k:q\in\mathcal B_K(k)}
  \Ree\left[
    \bm r_k\Hh\bm c_{kq}\frac{\partial s_q}{\partial z_i}
  \right]
  +\lambda_p
  \frac{\partial\mathcal R_{\rm code}}{\partial z_i},
  \label{eq:pfft_gradient}
\end{equation}
where the outer sum runs over the carriers that see carrier $q$.  A bit
therefore feels every carrier to which it contributes ICI, and the parity checks it takes part
in, at the same time.

\subsection{Combiner Symbol Tie}

A receiver that forms a combined symbol $\widehat s_k=\bm w_k\Hh\bm Y_k$ already
has an estimate of the transmitted symbol, and that estimate should agree with
the one the latents imply.  We add
\begin{equation}
  \mathcal L_{\rm tie}
  =
  \rho\sum_{k\in\mathcal D}\abs{\widehat s_k-s_k(\bm z)}^2 ,
\end{equation}
where $\rho$ sets how strongly the two estimates are tied together.  For bit $i$
on carrier $k$,
\begin{equation}
  \frac{\partial \mathcal L_{\rm tie}}{\partial z_i}
  =
  -2\rho\Ree\left[
    \big(\widehat s_k-s_k\big)^*
    \frac{\partial s_k}{\partial z_i}
  \right],
  \label{eq:tie_gradient}
\end{equation}
where the sign is negative because the residual is written as
$\widehat s_k-s_k$, and $s_k$ is the variable being differentiated.

\subsection{Sparse or Clustered Channel Refinement}

A sparse or clustered channel is described by a few delays and Doppler shifts,
and those few numbers are continuous.  After the support has been selected, we
write the model as
\begin{equation}
  \widehat{\bm y}=A(\bm\theta)\bm s(\bm z),\qquad
  \bm e=\widehat{\bm y}-\bm y,
\end{equation}
where $\bm\theta$ collects the continuous path or cluster parameters.  For a
scalar parameter $\theta_c$,
\begin{equation}
  \frac{\partial}{\partial\theta_c}\norm{\bm e}^2
  =
  2\Ree\left[
    \bm e\Hh
    \frac{\partial A(\bm\theta)}{\partial\theta_c}
    \bm s(\bm z)
  \right],
  \label{eq:theta_gradient}
\end{equation}
where $\partial A/\partial\theta_c$ is the sensitivity of the mixing matrix to
that parameter.  We use this for a small number of refinements only, such as a
residual Doppler offset or a cluster delay drift.  However, interpolation
boundaries make $\partial A/\partial\theta_c$ discontinuous for some
dictionaries, and a local grid search is then the better estimator.

\subsection{Impulse and External-Interference Nuisance Terms}

Snapping shrimp noise and partial-band interference both add energy that the
signal model cannot explain, and both are sparse in a domain the receiver knows.
We write
\begin{equation}
  \widehat{\bm y}=A\bm s(\bm z)+B\bm u,\qquad
  \bm e=\widehat{\bm y}-\bm y,
\end{equation}
where $B$ maps the nuisance terms into the observation.  Given $\bm z$, we
update $\bm u$ by
\begin{equation}
  \bm u
  =
  \argmin_{\bm u}\norm{A\bm s(\bm z)+B\bm u-\bm y}^2
  +\lambda_u\mathcal R_u(\bm u),
\end{equation}
which is least squares when $\mathcal R_u=0$, and a sparse proximal step when
$\mathcal R_u$ is an $\ell_1$ penalty.  Given $\bm u$, we update $\bm z$ by
\eqref{eq:universal_z_gradient}.  Alternating between the two is more reliable
than differentiating through a nuisance detector, because each update is the
exact solution of its own subproblem.

\section{The Solver Schedule}

We run the receiver as follows.
\begin{enumerate}
  \item Generate candidates for coarse Doppler, the number of partial-FFT views $M$, and
  ICI span $K_{\rm ICI}$.
  \item For each candidate, initialize the channel and response coefficients from
  the pilots, with sparse support selection where it applies.
  \item Alternate:
  \begin{enumerate}
    \item update the channel and interference variables by least squares,
    grid search, or a proximal rule;
    \item update $\bm z$ by a few hand-gradient steps using
    \eqref{eq:full_z_gradient};
    \item run BP, or use BP extrinsic information, to refresh the
    log-likelihood ratios and the confidence.
  \end{enumerate}
  \item Score the candidates by pilot residual, data residual, parity residual,
  log-likelihood ratio margin, BP success, and CRC.
  \item Report a high parity residual or a failed CRC as an erasure to the outer
  erasure code or to the adaptive modulation and coding policy.
\end{enumerate}
Steps 1 and 5 are where the discrete decisions are made, and neither of them
uses a gradient.

\section{What Not to Differentiate}

A hand-derived gradient is useful when the mapping is smooth and
low dimensional, and it is not the right tool for every part of these receivers.
We do not differentiate the following.
\begin{itemize}
  \item Integer quantities: the ICI span, the number of partial-FFT views, the pilot pattern, the
  LDPC graph, the support size, and the modulation mode.  A gradient with
  respect to an integer is not defined, and rounding a continuous relaxation
  gives a receiver that neither setting describes.
  \item Hard decisions: OMP support selection, GLRT decisions, CRC decisions,
  and BP stopping rules.  These are already the exact answer to their own
  subproblem.
  \item The finite-field LDPC algebra.  We keep BP explicit and use the relaxed
  parity residual only as a soft constraint at the receiver, so that the code
  remains the code that the transmitter used.
  \item The nuisance terms alone.  Impulse and interference variables can
  explain any observation if they are given enough freedom, so we always score a
  candidate by parity and CRC, and the receiver then prefers explanations that
  are also valid codewords.
\end{itemize}

\section{Conclusions}

We considered which parts of an underwater acoustic OFDM receiver benefit from a
gradient.  The coded-bit latents do, because the loss is smooth in them and they
are the variables the parity checks constrain.  Most of the physics does not,
because it is either discrete or linear, and is better estimated by grid search
or least squares.  The receiver that follows is not fully differentiable.  It
alternates least squares, grid search and proximal steps for the physics, a
hand-derived gradient for the coded bits, and BP with parity and CRC for code
consistency.  Partial FFT, progressive ICI, Doppler compensation, sparse and
clustered channels, interference cancellation, multiuser reception, and
physical-layer network coding all fit this pattern.  They change the observation
matrix or nuisance basis, while the coded-bit gradient of
\eqref{eq:universal_z_gradient} stays as it is.

\end{document}
